\documentclass[11pt,a4paper]{article}

% -----------------------------------------------------------------------------
% Packages and document settings
% -----------------------------------------------------------------------------
\usepackage[T1]{fontenc}
\usepackage[utf8]{inputenc}
\usepackage[english]{babel}
\usepackage{lmodern}
\usepackage{microtype}
\usepackage[a4paper,margin=2.5cm]{geometry}
\usepackage{setspace}
\usepackage{graphicx}
\usepackage{booktabs}
\usepackage{array}
\usepackage{threeparttable}
\usepackage{caption}
\usepackage{subcaption}
\usepackage{float}
\usepackage{siunitx}
\usepackage{amsmath}
\usepackage{amssymb}
\usepackage{xcolor}
\usepackage{natbib}
\usepackage[hidelinks]{hyperref}
\usepackage{fancyhdr}

\graphicspath{{../../outputs/paper/figures/}}

\sisetup{
  round-mode=places,
  round-precision=3,
  table-number-alignment=center
}

\captionsetup{
  font=small,
  labelfont=bf,
  justification=centering
}

\setstretch{1.15}
\setlength{\parindent}{0pt}
\setlength{\parskip}{0.6em}

\pagestyle{fancy}
\fancyhf{}
\fancyhead[L]{Corporate Bankruptcy Risk Prediction}
\fancyhead[R]{Machine Learning for Finance}
\fancyfoot[C]{\thepage}
\setlength{\headheight}{14pt}

\hypersetup{
  pdftitle={Corporate Bankruptcy Risk Prediction with Machine Learning},
  pdfauthor={Berke Pehlivan},
  pdfsubject={Machine Learning for Finance research paper}
}

% -----------------------------------------------------------------------------
% Title information
% -----------------------------------------------------------------------------
\title{\textbf{Corporate Bankruptcy Risk Prediction\\with Machine Learning}\\[0.6em]
\large Comparing Interpretable Logistic Regression and Tree-Based Models}
\author{Berke Pehlivan\\
\small Machine Learning for Finance\\
\small University of Bonn}
\date{Academic Year 2025/2026}

\begin{document}

\begin{titlepage}
  \centering
  \vspace*{2.5cm}
  {\LARGE\bfseries Corporate Bankruptcy Risk Prediction\\[0.25em]
  with Machine Learning\par}
  \vspace{0.8cm}
  {\large Comparing Interpretable Logistic Regression\\
  and Tree-Based Models\par}
  \vfill
  {\large Berke Pehlivan\par}
  \vspace{0.3cm}
  {\normalsize Machine Learning for Finance\par}
  {\normalsize University of Bonn\par}
  {\normalsize Academic Year 2025/2026\par}
  \vfill
\end{titlepage}

\pagenumbering{roman}

\begin{abstract}
% Write the abstract after all substantive sections are complete.
\end{abstract}

\tableofcontents
\clearpage
\pagenumbering{arabic}

% -----------------------------------------------------------------------------
% Main paper
% -----------------------------------------------------------------------------
\section{Introduction}
\label{sec:introduction}

\subsection{Financial motivation}

Corporate bankruptcy is an important problem for lenders, investors, auditors,
and risk managers because a firm's failure can impose losses on creditors and
shareholders and can disrupt relationships with employees, suppliers, and other
firms. Financial statements provide a structured record of a company's assets,
liabilities, revenues, expenses, and profitability. These variables may contain
signals of financial distress before failure occurs, which makes bankruptcy
prediction a natural application of statistical learning in finance. Early
research established that combinations of accounting variables can help
distinguish distressed from healthy firms \citep{altman1968}, while later work
showed how bankruptcy risk can be studied in a probabilistic classification
framework \citep{ohlson1980}.

Despite this long research history, bankruptcy prediction remains difficult.
Failure is a rare outcome, so a model can achieve high overall accuracy by
predicting that almost every observation remains alive. Such a model may appear
successful while failing to identify the firms that matter most for a distress
screening exercise. In addition, financial variables are often strongly
correlated, measured on very different scales, and affected by firm size.
Relationships between accounting variables and failure may also be nonlinear or
depend on interactions that are not represented by a simple linear model.
Finally, false negatives and false positives do not generally have the same
economic cost: missing a failed firm may be costly, but flagging too many healthy
firms can make a screening system impractical.

These features make model evaluation as important as model fitting. Accuracy is
reported for completeness, but it cannot be the main criterion in an imbalanced
setting. This study therefore emphasizes precision and recall for the failed
class, the failed-class F1-score, balanced accuracy, the area under the receiver
operating characteristic curve (ROC-AUC), and the area under the
precision--recall curve (PR-AUC). Precision--recall analysis is especially
informative when the negative class greatly outnumbers the positive class,
because it makes the false-positive burden visible rather than allowing the many
true negatives to dominate the assessment \citep{davis2006}. The probability
threshold is also treated as a decision choice rather than assuming that the
default value of 0.50 is economically optimal.

\subsection{Research question and contribution}

This paper investigates the following research question:

\begin{quote}
\emph{How well can company-year financial statement variables identify failed
observations out of sample, and do flexible tree-based classifiers improve on an
interpretable Logistic Regression benchmark?}
\end{quote}

The analysis uses the American corporate bankruptcy dataset introduced by
\citet{lombardo2022}. The working sample contains 78,682 company-year
observations from 8,971 anonymized companies between 1999 and 2018, with 18
financial predictors and a binary failed/alive target. Because companies may
appear in multiple years, observations are split at the company level. This
prevents the same company from appearing in both training and final-test data
and reduces the risk that persistent company characteristics create an
artificially favorable evaluation. The design tests generalization to unseen
companies; it is not presented as a chronological forecasting backtest.

The empirical comparison begins with a majority-class baseline, which
demonstrates why accuracy alone is misleading. Logistic Regression provides the
main interpretable benchmark, and L1- and L2-regularized variants examine whether
coefficient shrinkage improves generalization when predictors are correlated.
A Decision Tree, Random Forest, and Gradient Boosting model then test whether
nonlinear rules and interactions improve predictive performance. These models
represent a progression from a transparent parametric classifier to more
flexible ensemble methods, reflecting the interpretability--flexibility and
bias--variance trade-offs emphasized in statistical learning
\citep{james2023}. A PCA plus Logistic Regression specification is included as a
dimension-reduction extension, not as the main model.

The paper makes three practical contributions within the scope of a course
project. First, it constructs a reproducible and leakage-aware comparison of
interpretable and tree-based classifiers. Second, it evaluates the models using
metrics that focus directly on the minority failed class and uses an untouched
final test set for the principal comparison. Third, it connects statistical
performance to the financial screening problem by examining confusion-matrix
counts, validation-based threshold trade-offs, Logistic Regression
coefficients, and tree-based feature importance. The resulting estimates are
interpreted as predictive associations rather than causal effects. The study is
therefore an empirical model comparison and not a claim that the fitted models
are ready for automated credit or investment decisions.

\section{Data}
\label{sec:data}

\subsection{Dataset and unit of observation}

The project uses the American corporate bankruptcy dataset published by
\citet{lombardo2022}. It contains accounting information for public companies
listed in the United States. The version used in this project has 78,682 rows
covering the years 1999 to 2018. It contains 8,971 anonymized company
identifiers, 18 financial predictors, a year variable, and a company status.
The data summary reports no missing values and no duplicated rows.

One row represents one company in one year. The same company can therefore
appear in several years. This panel-like structure is important for the way the
sample is split. A random split of individual rows could place one company's
earlier observations in the training sample and its later observations in the
test sample. The model might then benefit from company-specific information that
is already present in training. To reduce this risk, the project assigns whole
companies to the training or final-test sample. The split design is explained in
Section~\ref{sec:methodology}.

The source data cover a long period in which both the set of listed companies
and the wider economic environment changed. However, the project does not use a
time-based holdout. The final test set contains unseen companies, but its years
overlap with the years in the training set. The results should therefore be read
as evidence about generalization across companies under this sample design, not
as a real-time forecast of a later economic period.

\subsection{Target and financial predictors}

The original status is recorded as \texttt{alive} or \texttt{failed}. For
modeling, failed observations are coded as 1 and alive observations as 0. The
failed class is treated as the positive class throughout the analysis. The
company identifier and year are kept for splitting and reporting, but they are
not used as predictors. This leaves 18 financial inputs, labelled \texttt{X1}
to \texttt{X18} in the source file.

The predictors cover several parts of a company's financial position and
performance. They include current assets, inventory, receivables, market value,
and total assets; profitability measures such as EBITDA, EBIT, net income,
gross profit, and retained earnings; long-term debt, current liabilities, and
total liabilities; and activity measures such as sales, revenue, cost of goods
sold, and operating expenses. Depreciation and amortization is also included.
The model therefore uses reported financial levels rather than a set of
engineered financial ratios. Readable names are retained in the tables and
figures so that the model outputs can be discussed in financial terms.

These variables have very different scales and several describe related parts
of the same financial statements. For example, revenue and net sales are closely
connected, as are assets and liabilities. This matters for Logistic Regression,
because coefficient size depends on measurement scale and correlated predictors
can make individual coefficients less stable. The Logistic Regression and PCA
pipelines therefore standardize the predictors using parameters learned from
the training data. Tree-based models do not require this scaling.

\subsection{Class imbalance and descriptive evidence}

The target is strongly imbalanced. Of the 78,682 company-year observations,
73,462 are alive and 5,220 are failed. The failed class therefore represents
only 6.63\% of the full sample. Figure~\ref{fig:class-distribution} shows this
difference. A classifier that predicts every observation as alive would be
correct in more than 93\% of the full sample, although it would detect no
failures. This is the main reason why the later evaluation does not rely on
accuracy alone.

\begin{figure}[H]
  \centering
  \includegraphics[width=0.72\textwidth]{class_balance.png}
  \caption{Class distribution in the full dataset. The failed class accounts
  for 6.6\% of company-year observations.}
  \label{fig:class-distribution}
\end{figure}

The share of failed observations also changes over time. As shown in
Figure~\ref{fig:annual-failure-rate}, the observed rate rises from 7.16\% in
1999 to a peak of 9.40\% in 2003 and then generally declines, reaching 1.32\%
in 2018. This pattern should be interpreted with care. It may reflect changes in
economic conditions, but it may also reflect changes in the companies included
in the dataset or in the way the status is represented. The figure is used as a
description of the sample and not as evidence of a causal time trend.

\begin{figure}[H]
  \centering
  \includegraphics[width=0.82\textwidth]{annual_failure_rate.png}
  \caption{Observed annual failure rate in the full dataset, 1999--2018.}
  \label{fig:annual-failure-rate}
\end{figure}

Table~\ref{tab:financial-medians} compares the medians of six selected
financial variables. Medians are used because financial statement values are
often skewed and can contain very large observations. Failed observations have
lower median market value, net income, current assets, and retained earnings.
They have higher median long-term debt and total liabilities. The signs are
financially plausible as descriptive patterns, but they do not show that these
variables cause failure. They are simple comparisons made before controlling
for the other predictors.

\begin{table}[H]
  \centering
  \caption{Selected financial-variable medians by company status}
  \label{tab:financial-medians}
  \begin{threeparttable}
    \begin{tabular}{lrr}
      \toprule
      Financial variable & Alive median & Failed median \\
      \midrule
      Market value         & 240.90 & 117.80 \\
      Net income           &   2.07 &  -3.33 \\
      Total long-term debt &   7.09 &  18.44 \\
      Current assets       & 102.92 &  75.87 \\
      Total liabilities    &  80.74 &  97.04 \\
      Retained earnings    &  -0.17 & -25.51 \\
      \bottomrule
    \end{tabular}
    \begin{tablenotes}[flushleft]
      \footnotesize
      \item Notes: Values are shown in the units reported by the source
      dataset. The table is descriptive and does not imply causal effects.
    \end{tablenotes}
  \end{threeparttable}
\end{table}

\section{Methodology}
\label{sec:methodology}

\subsection{Company-level training, validation, and test design}

\subsection{Majority-class baseline}

\subsection{Logistic Regression and regularization}

\subsection{Tree-based models}

\subsection{PCA extension}

\section{Evaluation Strategy}
\label{sec:evaluation}

\subsection{Why accuracy is insufficient}

\subsection{Ranking and failed-class metrics}

\subsection{Threshold analysis}

\section{Results}
\label{sec:results}

\subsection{Final-test model comparison}

\subsection{Classification trade-offs}

\subsection{Precision--recall analysis}

\section{Model Interpretation}
\label{sec:interpretation}

\subsection{Logistic Regression coefficients}

\subsection{Tree-based feature importance}

\section{PCA Results}
\label{sec:pca-results}

\section{Critical Reflection and Limitations}
\label{sec:limitations}

\section{Conclusion}
\label{sec:conclusion}

% -----------------------------------------------------------------------------
% References
% -----------------------------------------------------------------------------
\clearpage
\bibliographystyle{apalike}
\IfFileExists{references.bib}
  {\bibliography{references}}
  {\section*{References}}

% -----------------------------------------------------------------------------
% Appendix
% -----------------------------------------------------------------------------
\clearpage
\appendix

\section{Additional Model Results}
\label{app:model-results}

\section{Threshold Analysis}
\label{app:thresholds}

\section{PCA Diagnostics}
\label{app:pca}

\section{Reproducibility and Code Structure}
\label{app:reproducibility}

\end{document}
